% Book-specific commands, environments and metadata.

% Title metadata used by the frontmatter pages.
\newcommand{\thebooktitle}{Splitting the Graph}
\newcommand{\thebooksubtitle}{Federated GraphQL with Hot Chocolate and Cosmo on .NET}
\newcommand{\thebookauthor}{Giang Dang}

% Inline code in prose.
\newcommand{\code}[1]{\texttt{#1}}

% Diaeresis on a cited author's name. SPEC decision 30 puts the book in strict
% Ascii mode, so a name like Philip Mueller's is spelled with an accent macro
% rather than the Unicode letter; spelling it "Mueller" instead would misspell
% a real person. This wrapper exists because writing \" directly in a chapter
% trips the prose gate's quote check, which reads the accent's double quote as
% a literal quotation mark. Usage: M\uml{u}ller.
\newcommand{\uml}[1]{\"{#1}}

% Acute accent, the same decision as \uml one letter along. Chapter 12 cites
% Tom Houle, whose name carries one, and the reasoning is identical: strict
% Ascii mode forbids the Unicode letter, respelling him would misspell a real
% person, and a bare \' in a chapter file is an apostrophe as far as the prose
% gate's quote check is concerned. Usage: Houl\acu{e}.
\newcommand{\acu}[1]{\'{#1}}

% The one skin every box in this book wears, so that a callout and a summary
% read as the same family: square corners, a single 0.4pt stroke, black on
% white, and a sans-serif title in its own compartment above a rule. That is
% the box equivalent of the figure idiom in SPEC decision 33.
%
% Two settings here are load-bearing and were arrived at by building the box
% and looking at the page, so do not drop them:
%
%   coltitle=black    tcolorbox derives the title color from colframe, which
%                     is black, so the default title color is white. On a
%                     white title background that renders an invisible title.
%
%   the title compartment, rather than "attach boxed title to top left"
%                     the attached-tab form draws its tab but then draws only
%                     the left rule of the box itself: no top, right or bottom
%                     edge. The compartment form draws all four.
%
% Both defects were present in the first version of the onhc14 box below and
% are what a summary box built to copy it inherited.
\tcbset{
  bookbox/.style={
    enhanced,
    breakable,
    colback=white,
    colbacktitle=white,
    colframe=black,
    coltitle=black,
    boxrule=0.4pt,
    titlerule=0.4pt,
    arc=0pt,
    outer arc=0pt,
    left=3mm, right=3mm, top=2mm, bottom=2mm,
    fonttitle=\sffamily\bfseries\small,
  }
}

% SPEC decision 32: the "On HotChocolate 14" callout.
%
% Hot Chocolate 16 is this book's spine and 14 is out of support. A reader on
% 16 must be able to skip every one of these without losing the thread, which
% is why they are boxed rather than set as body prose. A callout may only
% state what the verification repo's hc14 branch actually compiled (decision
% 11), and it points at appendix B for the full mapping.
%
% Usage:
%   \begin{onhc14}
%     ...prose, and listings if the branch compiled them...
%   \end{onhc14}
\newtcolorbox{onhc14}{bookbox, title={On HotChocolate 14}}

% SPEC decisions 20 and 21: the end-of-chapter summary box.
%
% This book has no labs. Every chapter closes with one of these instead, and
% the box is the single named exception to the Economy line, which is why its
% shape is fixed rather than left to whoever drafts the chapter.
%
% Decision 21 fixes two shapes. Parts I, II and IV close with three to five
% bullets stating the chapter's claims:
%
%   \begin{chaptersummary}
%     \item ...
%   \end{chaptersummary}
%
% Part III closes with a fixed symptom, cause, fix triple - the symptom a
% reader searches for months later, the mechanism underneath it, and the change
% that fixes it. That triple is what makes the box a lookup rather than a
% restatement, so it takes three mandatory arguments and cannot be written any
% other way:
%
%   \begin{problemsummary}{<symptom>}{<cause>}{<fix>}\end{problemsummary}
%
% Both wear the shared bookbox skin above, so a summary and a callout are
% visibly the same object.
\newtcolorbox{summarybox}[1]{bookbox, title={#1}}

\newenvironment{chaptersummary}
  {\begin{summarybox}{Summary}\begin{itemize}[leftmargin=*, topsep=0pt, itemsep=2pt]}
  {\end{itemize}\end{summarybox}}

\newenvironment{problemsummary}[3]
  {\begin{summarybox}{Summary}%
   \begin{description}[leftmargin=0pt, style=sameline, font=\sffamily\bfseries, topsep=0pt, itemsep=2pt]
   \item[Symptom.] #1
   \item[Cause.] #2
   \item[Fix.] #3
   \end{description}}
  {\end{summarybox}}

% Theorem environments (amsthm) - uncomment if the book needs them.
% \theoremstyle{definition}
% \newtheorem{definition}{Definition}[chapter]
